\subsection{2026-09-21 --- Centralizing the scattered data into a modeling table}
\label{sec:modeling-table}

\subsubsection*{Motivation.} The raw source files are scattered across several CSVs at different
native granularities (hourly price/inflow, daily storage/water value, one row per 168-hour case
for labels). Keeping them separate is fine, and good practice, as the unmodified source of
truth. But a supervised model needs one row = one example, features..., target at the end ---
so before any modeling, everything is joined once into a flat table, which also forces explicit,
one-time decisions about how sources at different granularities align to a row and where the
train/test boundary falls.

\subsubsection*{Design decisions (discussed and agreed before building).}
\begin{itemize}
  \item \textbf{Row granularity: one row per (case, hour).} Considered alternatives: one row per
    (case, hour, unit) (long format, $\sim$6.9M rows, simplest per-row model but duplicates
    case/hour features 14$\times$ and drops the multi-unit correlation structure from
    Section~\ref{sec:eda}), and one row per case (fully wide, 2{,}922 rows but 2{,}352 target
    columns, awkward for standard tabular models). Case+hour keeps a manageable row count while
    keeping all 14 units' targets on one row, so a model can see them jointly if we want it to.
  \item \textbf{Raw joins only, no engineered features.} No normalization, no lags, no derived
    columns at this stage --- feature engineering (e.g.\ the within-case price normalization
    from Section~\ref{sec:eda}) is treated as a separate, explicit step.
  \item \textbf{Core dataset (\texttt{data/kernel/}) only.} Extended data (topology,
    constraints) intentionally left out until there's a concrete reason to use it.
  \item \textbf{Per-reservoir columns kept separate}, not pre-aggregated (17 inflow + 17 storage
    + 17 water-value columns, individually named), so nothing is thrown away before we've
    decided how to use it.
\end{itemize}

\subsubsection*{Implementation.} \texttt{src/build\_modeling\_table.py}. Daily-native features
(storage, water value) are repeated across each day's 24 hours since they're only observed once
per day. Targets are reshaped from the labels file's wide format
(\texttt{result\_committed\_<unit>\_t<hour>}) into one \texttt{on\_<unit>} column per unit.
Output:

\begin{itemize}
  \item \texttt{data/processed/train\_hourly.parquet} --- 490{,}896 rows (2{,}922 cases
    $\times$ 168 hours) $\times$ 72 columns (5 id/time columns, 53 raw features, 14 targets).
  \item \texttt{data/processed/test\_hourly.parquet} --- 121{,}800 rows (725 cases $\times$
    168 hours) $\times$ 58 columns (same features, no targets).
\end{itemize}

\finding{Spot-checked the join against raw source files (a mid-case UC label, a price value, a
storage value at a specific timestamp) --- all matched exactly. No missing values after joining
any feature column.}
